\chapter{Sampling, Spectra, and the Limits of a Reading}
\label{chap:sampling-spectra-limits}

Every reading the instrument takes is a projection. A continuous field, infinite in its detail, is mapped
onto a finite ledger of counts, and a projection always does two things at once: it keeps part of what it is
given and discards the rest. The part it keeps is the reading; the part it discards is its \emph{kernel} ---
the blind spot every projection has, the subspace of the world it is structurally unable to see. This chapter
is a tour of that blind spot. It asks when a reading is dense enough to lose nothing, what a finite reading
does to a sharp edge, where the edge goes when the reading smooths it away, what two finite readings invent
between them that neither holds alone, and what part of a message no receiver can ever decode. Five effects,
one object: the kernel of the projection, and where in physics it hides.

\section{The Fourier--Nyquist Effect: exact inversion needs density}
\label{se:fourier-nyquist}

The instrument can take a continuous wave apart into its frequencies and put it back together exactly --- but
only if it sampled densely enough to begin with. Taking the wave apart is the Fourier transform, and putting
it back is its inverse,
\begin{equation}
  \hat f(k) = \int f(x)\, e^{-ikx}\, dx, \qquad f(x) = \frac{1}{2\pi}\int \hat f(k)\, e^{ikx}\, dk,
\end{equation}
the forward integral reading off how much of each frequency $k$ the signal holds, the inverse rebuilding the
signal from that spectrum; the two are exact inverses, so the spectrum $\hat f$ carries everything the signal
does, written in frequencies instead of positions. The sampling enters when the instrument records the signal
not as a continuous curve but as a row of samples. A signal is band-limited when its spectrum stops at a
highest frequency $B$,
\begin{equation}
  \hat f(k) = 0 \quad\text{for } |k| > B,
\end{equation}
holding no detail finer than $B$; and Nyquist found, Shannon made a theorem, that such a signal is recovered
exactly from its samples if and only if they come at least twice that fast,
\begin{equation}
  f_s \geq 2B .
\end{equation}
Sample at or above the Nyquist rate and the reconstruction is perfect, the discrete reading carrying every
distinction the continuous wave held. Sample below it, and high frequencies fold down onto low ones ---
aliasing --- so that two genuinely different waves leave exactly the same trail of samples. The under-sampled
reading cannot tell them apart, not because noise has blurred them but because the projection has
\emph{identified} them: they differ only in the kernel.

Everyone has seen the kernel at work. Film a spinning wagon wheel at too few frames per turn and the wheel
appears to rotate slowly backward, or to hang motionless --- the camera, sampling below the wheel's Nyquist
rate, has filed a fast forward spin and a slow backward one as the same record. The true motion has not been
corrupted; it has been folded onto another that shares its samples, and from the film alone the two are one.
That folding is the kernel of an under-dense reading: the set of distinct worlds a too-sparse sampling cannot
separate.

The honest floor is a smooth-shadow analogy. The finite sampling is the discrete shadow of the continuum
signal, and exact inversion is the bridge between them --- a bridge that holds above the Nyquist rate and
collapses below it, where the hinge of the analogy is precisely the density of the sampling. The reading is
that a reading recovers everything if and only if it is dense enough. Below the Nyquist rate, distinct signals
collapse to identical samples, and the loss is not noise added but the geometry of under-sampling: two
different worlds filed as one, separated only by what the reading threw away.

\section{The Gibbs Phenomenon: the overshoot that won't go away}
\label{se:gibbs}

Push a sharp edge through a band-limited reading and the reading rings. At a discontinuity --- a square wave's
vertical face, a brightness step in an image --- a band-limited reconstruction overshoots the jump and then
oscillates beside it, and the overshoot does not shrink as more frequencies are added: it stays, fixed, at
about nine percent of the jump's height. Add a hundred terms, a thousand, a million, and the ripple narrows
but its peak holds at the same nine percent. The number is exact and famous. A partial Fourier sum --- the
first $N$ terms of the series --- reconstructs a jump by convolving it with the Dirichlet kernel, the summed
first $N$ harmonics, and that kernel has a fixed overshoot at a step: as $N$ grows the ripple is squeezed
toward the edge but its first peak settles at a constant fraction of the jump, the Gibbs constant, about
$0.0895$ --- nine percent --- no matter how many terms are summed. The overshoot does not converge away because
it is a property of the kernel, not of how finely the sum is taken. In the ledger's terms the part of a jump of
height $h$ that a band-limit $B$ cannot represent is
\begin{equation}
  \mathrm{gap} = h - \min(h, B),
\end{equation}
and that gap is what the ringing is made of.

The ripple is familiar wherever sharp meets smooth. Reconstruct a square wave from a truncated Fourier series
and the telltale ears appear at every corner; compress an image too hard and a faint halo rings each crisp
boundary; the overshoot is the same in each case, and in each case it cannot be filtered away without blurring
the edge that causes it. The honest floor is a smooth-shadow analogy: the overshoot is the discrete jump's
smooth shadow, and it persists precisely because the shadow cannot represent the discontinuity it shadows. The
reading is that a band-limited reading cannot show a sharp edge without ringing. The nine percent is not a
flaw in the signal and not a fault in the instrument; it is the fixed price of the shadow, the irreducible
residue a continuous reading pays whenever it is asked to draw a corner it cannot hold.

\section{The Gibbs Preservation Effect: edges in the null space}
\label{se:gibbs-preservation}

If the previous effect showed the kernel as a price, this one shows it as a refuge, and the two are the same
kernel seen twice. The feature a band-limited reading cannot represent --- the sharp edge that makes it ring
--- is the very feature a smoothing cannot reach to erase. The edge that overshoots in the one reading is the
edge preserved in the other; the null space that rings is the null space that keeps. The instrument's blind
spot exacts a price at a corner and offers a refuge at the same corner, because it is one blind spot met from
two sides --- and seeing that is the clearest statement of the chapter's whole idea.

When the instrument smooths a record --- projects a discrete ledger onto its continuous shadow --- it applies a
band-limiting projection $P$, a map that keeps the low frequencies and zeroes the rest. Like every projection,
$P$ splits its domain in two: the range, the smooth part it keeps, and the null space, the features it sends to
nothing, $P e = 0$. A sharp antisymmetric feature --- an edge of the form $[+h, -h]$ --- lies squarely in that
null space: $P$ maps it to zero. To the smoothing the edge is invisible. Add it to the record or remove it, and
the smooth shadow $P f$ does not change, because the edge sits entirely in the part the projection discards,
never in the part it keeps.

The consequence is the opposite of what intuition expects. One imagines smoothing as the enemy of edges --- a
blur that erodes them. But an edge that lives in the kernel of the blur is not eroded; it is \emph{set aside}.
Blur a photograph with a sharp boundary and the boundary is not smeared into the soft regions around it; it is
separated, intact, into the residual the blur discards, and from that residual it can be recovered whole. The
smoothing cannot destroy the edge for exactly the reason it cannot see it: the edge was never in the
smoothing's reach to begin with. The instrument's blind spot, here, is also its safe-deposit box --- the very
subspace the projection cannot read is the one place a feature is safe from it.

The honest floor is a smooth-shadow analogy --- the preserved edge is the part of the discrete ledger the
smooth shadow cannot touch, the analogy's blind spot named for what it is. The reading is that smoothing does
not destroy edges; it cannot reach them. The sharpest features of a record are precisely the ones the smooth
shadow leaves alone, which is why a reading can blur a field and still carry its discontinuities perfectly
intact: the edge was never in the shadow's range to lose, but in the kernel the shadow sets aside. This is the
kernel at its most vivid --- a blind spot that, far from losing what falls into it, keeps it.

\section{The Moir\'e Effect: beats from mismatched lattices}
\label{se:moire}

Lay two fine gratings over each other, their pitches slightly different, and a coarse pattern swims into view
--- broad bright and dark bands, far larger than either grating's lines, marching at the difference of the two
pitches,
\begin{equation}
  f_{\mathrm{beat}} = |f_1 - f_2| .
\end{equation}
The beat is the difference frequency two periodic patterns make when they are laid together. Overlay two
gratings of frequencies $f_1$ and $f_2$ and their product splits, by the elementary product-to-sum identity,
into a fast term at the sum $f_1 + f_2$ and a slow term at the difference $|f_1 - f_2|$; the eye averages away
the fast term and sees only the slow one, the coarse band marching at the difference. This is the moir\'e
effect, named for the watered silk whose rippled sheen comes from exactly such an overlay,
and seen daily wherever two regular patterns cross: a striped shirt shimmering on a television screen, the
rosettes that betray a halftone reprint, the bands across a photograph of a brick wall. The coarse pattern is
real --- one can measure it, photograph it, count its bands --- and yet it is in neither grating alone. It is
a structure the two readings make \emph{together}, a beat of their mismatch, present in the pair and absent
from each.

The honest floor is a finite count obligation: the beat is a definite difference of two pitches, the count of
where the two lattices fall into step against where they drift apart. The reading is that the moir\'e pattern
belongs to the pair and not to the world. It is what two finite readings invent between them --- a ghost of
their mismatch, read as a structure that was never in the signal either reading sampled. Where the kernel of a
single reading is what that reading cannot see, the moir\'e pattern is its mirror image: a thing two readings
together can see that is not, in fact, there.

\section{The Message Effect: what no receiver can decode}
\label{se:message}

A transmission carries structure, but a receiver registers only the part its own basis can resolve. Relative
to that basis, the symmetric component of a refinement is decodable, and the antisymmetric twist is not: it
falls in the kernel of the decoding, and for that receiver it is simply unrecordable. Only the symmetric part
survives as message; a pure antisymmetric twist arrives reading identical to silence, indistinguishable from
nothing sent at all. The unrecoverable part has a precise signature --- it is a closure defect, a circulation
that does not vanish around a loop,
\begin{equation}
  \oint \mathbf{R}\cdot d\boldsymbol{\ell} \neq 0,
\end{equation}
so that the twist which refuses to close around a loop is exactly the twist that cannot be sent.

This is the kernel as a limit on communication, and it is sharper than it first sounds. It does not say the
twist is hard to decode, or that a better receiver could recover it; it says that \emph{for this basis} the
twist and silence are the same message. Two transmissions differing only by a pure twist are, to the receiver,
one transmission. The honest floor is a smooth-shadow analogy --- the decodable message is the symmetric
shadow of the transmission, and the antisymmetric twist is the part the projection cannot cast, named as the
kernel of the decoding. The reading is that communication is a projection, and a projection has a kernel. What
no receiver can decode is not lost in transit but invisible to the basis --- a twist that reads as silence
because the decoder has no axis for it. The limit of a reading is its kernel, and a message is only ever its
symmetric part.

\bigskip

With the message, the chapter's object is complete, and it has a single name. A finite reading is a
projection; a projection keeps and discards; and what it discards --- its kernel --- is the blind spot every
reading carries. The kernel appears here as a threshold below which two worlds file as one, as the ringing a
shadow cannot smooth away, as the refuge where an edge survives precisely because the smoothing cannot reach
it, as the ghost two mismatched readings invent together, and as the twist that reaches every receiver reading
exactly like silence. The instrument is honest about all of it: it names what each reading recovers and what
it leaves behind, and never mistakes the second for nothing. Part II closes next with the inverse-square law
--- a field's reach diluting across an expanding frontier --- and with it the instrument leaves the reading of
waves for the geometry of how a field spreads.
